\documentclass[12pt,a4paper]{article}

\usepackage[utf8]{inputenc}
\usepackage[T1]{fontenc}
\usepackage[french]{babel}

\usepackage{amsmath, amssymb}

\usepackage{graphicx}
\usepackage{caption}

\usepackage{siunitx} % Pour les unités

\usepackage{geometry}
\geometry{top=2cm, bottom=2.5cm, left=3cm, right=3cm}
\setlength{\parindent}{0pt}

\usepackage{enumitem}

\usepackage{hyperref}

\usepackage{xcolor}

\usepackage{multicol}
\setlength{\columnseprule}{1pt}
\def\columnseprulecolor{\color{gray}}

\usepackage{tcolorbox}


\usepackage{import}
\usepackage{xifthen}
\usepackage{pdfpages}
\usepackage{transparent}
\usepackage{float}

\newcommand{\incfig}[1]{%
	\def\svgwidth{\columnwidth}
	\import{./figures/}{#1.pdf_tex}
}


\newcommand{\RomanNum}[1] %pour les chiffre romain
{\MakeUppercase{\romannumeral #1}}











\title{Mon projet structure en BA V2}
\author{}
\date{\today}

\begin{document}
	\counterwithin{section}{part} % pour que \section{} recommence de 1 aprés \part{}
	
	
	\tableofcontents
	
	\newpage
	
	
	
	\part{Présentation et généralités}
	
	\section{Introduction}
	
	\section{Présentation de l'ouvrage}
	
	\section{Caractéristique des matériaux}
	
	
	
	
	\part{Prédimensionnement des éléments}
	
	Dans ce chapitre les éléments prédimensionné sont :
	
	\begin{itemize}
		\item Poteaux, poutres.
		\item Plancher.
		\item Voiles.
		\item Balcon.
		\item Escalier.
		\item Acrotére.
	\end{itemize}
	
	
	\section{Prédimensionnement du plancher}
	
	Il existe deux type de plancher en corps creux un 16(corps creux)+4(dalle de compression) et un 20+5.
	
	\[ H_{p} = \dfrac{L_{\max}}{22,5} \]
	
	Avec :
	
	$H_{p}$ : la hauteur du plancher.
	
	$L_{\max}$ : La plus grande travée (plus grande distance entre nue d'appuies).
	
	
	
	\[ \begin{cases}
		\text{si  } H_{p} < 16+4 = 20~\unit{cm} \text{  alors on choisie un plancher en 16+4} \\\\
		\text{sinon si  } 16+4 < H_{p} < 20+5 \text{  alors on va opter pour un plancher en 25 cm} 
	\end{cases} \]
	
	
	
	\[ H_{p} = \dfrac{440}{22,5} = 19,55 ~\unit{cm} < 20 ~\unit{cm} \]
	
	Donc un plancher de 16+4=20 cm.
	
	
	\section{Prédimensionnment des poutres}
	
	Comme le plancher est en corps creux donc on aura deux type de poutre des poutres principale et d'autre secondaire.
	
	\vspace{\baselineskip}
	
	Ce que impose le réglement parassismique algérien (RPA 99 version 2003).
	
	Dans l'article 7.5.1 on a pour le coffrage : 
	
	\begin{equation}
		\begin{aligned}
			b& \geq 20 ~\unit{cm} \\
			h& \geq 30 ~\unit{cm} \\
			\dfrac{h}{b}& \leq 4 \\
			b_{max}& \leq 1,5\cdot h + b
		\end{aligned}
		\label{RPA coffrage des poutres}
	\end{equation}
	
	
	\vspace{\baselineskip}
	
	\subsection{Hauteur}
	
	\[ \dfrac{L_{\max}}{15} \leq h \leq \dfrac{L_{\max}}{10} \]
	
	
	
	
	
	\[ \dfrac{440}{15} \leq h \leq \dfrac{440}{10} \]
	
	\[ 29,33 ~\unit{cm} \leq h \leq 44 ~\unit{cm} \]
	
	Finalement on opte pour :
	
	\[ h = 35 ~\unit{cm} \]
	
	\subsection{Largeur}
	
	\[ 0,4 \cdot h \leq b \leq 0,7 \cdot h \]
	
	\[ 14 ~\unit{cm} \leq b \leq 24,5 ~\unit{cm} \]
	
	On prend :
	
	\[ b = 20 ~\unit{cm} \]
	
	\[ \dfrac{h}{b} = \dfrac{35}{20} = 1,75 < 4 \] 
	
	{\centering Donc aucun risque de deversement. \par}
	
	\[ b_{\max} \leq 1,5 \cdot 35 + 20 = 72,5 ~\unit{cm} \]
	
	{\centering Condition vérifiée. \par} 	
	
	
	
	\vspace{\baselineskip}
	
	Dimensions retenue : 
	
	\[ \fbox{$h=35 ~\unit{cm} \quad ; \quad b=20 ~\unit{cm}$} \]
	
	
	
	
	
	
	
	
	
	
	
	
	
	
	
	
	
	
	
	
	
\end{document}